% One-time imported draft. Editable native source; no Markdown regeneration.
\section{Derive valid lookup
positions}\label{005-seven-segment-counter-x-window-bounds-derive-valid-lookup-positions}

For start f and count n, valid selections satisfy f \ensuremath{\geq} 0, n \ensuremath{\geq} 0, and f +
n \ensuremath{\leq} 16.

{\def\LTcaptype{table} % do not increment counter
\begin{longtable}[]{@{}
  >{\RaggedRight\arraybackslash}p{(\linewidth - 4\tabcolsep) * \real{0.3333}}
  >{\RaggedRight\arraybackslash}p{(\linewidth - 4\tabcolsep) * \real{0.3333}}
  >{\RaggedRight\arraybackslash}p{(\linewidth - 4\tabcolsep) * \real{0.3333}}@{}}
\toprule\noalign{}
\begin{minipage}[b]{\linewidth}\RaggedRight
Window
\end{minipage} & \begin{minipage}[b]{\linewidth}\RaggedRight
Selected indices
\end{minipage} & \begin{minipage}[b]{\linewidth}\RaggedRight
Cursor after completion
\end{minipage} \\
\midrule\noalign{}
\endhead
\bottomrule\noalign{}
\endlastfoot
f = 0, n = 16 & 0 through 15 & glyphs + 16 \\
f = 14, n = 2 & 14 and 15 & glyphs + 16 \\
f = 5, n = 1 & 5 & glyphs + 6 \\
f = 16, n = 0 & None & glyphs + 16 \\
\end{longtable}
}

The last row forms an exclusive-end pointer but does not dereference it
in the reference program. An empty request must not become a do-while
iteration.
